\documentclass[11pt,a4paper]{article}

% ── Codificación, fuentes y lenguaje ───────────────────────────────────────────
\usepackage[utf8]{inputenc}
\usepackage[T1]{fontenc}
\usepackage[default,scale=0.95]{sourcesanspro}
\renewcommand{\familydefault}{\sfdefault}
\usepackage[spanish,es-noshorthands,es-tabla]{babel}

% ── Geometría y espaciado ─────────────────────────────────────────────────────
\usepackage[top=2.2cm, bottom=2.5cm, left=2.2cm, right=2.2cm, headheight=15pt]{geometry}
\usepackage{setspace}
\setstretch{1.18}
\usepackage{parskip}

% ── Matemáticas y unidades ────────────────────────────────────────────────────
\usepackage{amsmath}
\usepackage{amssymb}
\usepackage{siunitx}

% ── Circuitos ─────────────────────────────────────────────────────────────────
\usepackage[european, straightvoltages]{circuitikz}

% ── Tablas ────────────────────────────────────────────────────────────────────
\usepackage{booktabs}
\usepackage{tabularx}
\usepackage{array}
\usepackage{longtable}
\usepackage{colortbl}
\usepackage{enumitem}

% ── Colores, cajas e iconos ───────────────────────────────────────────────────
\usepackage[dvipsnames]{xcolor}
\usepackage{tcolorbox}
\tcbuselibrary{skins, breakable}
\usepackage{fontawesome5}
\usepackage{graphicx}

% ── Listados de código (dark-mode infográfico) ────────────────────────────────
%   Estilo base: fondo oscuro con números de línea. Los bloques lstlisting se
%   envuelven automáticamente en una caja tcolorbox redondeada.
\usepackage{listings}

% ── Secciones con barra de color (infografía) ────────────────────────────────
\usepackage{titlesec}
\titleformat{\section}
  {\normalfont\LARGE\bfseries\color{azulNoche}}
  {\colorbox{cyanNeon}{\color{fondoOscuro}\thesection}}{0.7em}{}
  [\vspace{2pt}\color{cyanNeon}\rule{\linewidth}{1.8pt}]
\titlespacing*{\section}{0pt}{24pt}{10pt}
\titleformat{\subsection}
  {\normalfont\large\bfseries\color{azulNoche}}
  {\textcolor{cyanNeon}{\thesubsection}}{0.7em}{}
  [\vspace{1pt}\color{grisLinea}\rule{\linewidth}{0.6pt}]
\titlespacing*{\subsection}{0pt}{14pt}{6pt}
\titleformat{\subsubsection}
  {\normalfont\normalsize\bfseries\color{azulNoche}}
  {\textcolor{cyanNeon}{\thesubsubsection}}{0.7em}{}
\titlespacing*{\subsubsection}{0pt}{10pt}{4pt}

% ── Cabeceras y pies de página ────────────────────────────────────────────────
\usepackage{fancyhdr}
\usepackage{lastpage}
\pagestyle{fancy}
\fancyhf{}
\renewcommand{\headrulewidth}{0pt}
\renewcommand{\footrulewidth}{0pt}
\fancyfoot[C]{%
  \begin{minipage}{\linewidth}
    \vspace{2pt}
    \color{cyanNeon}\rule{\linewidth}{1.2pt}\\[2pt]
    \centering\color{grisTexto}\footnotesize
    Página \thepage\ de \pageref{LastPage}
  \end{minipage}%
}

% ── Hiperenlaces ──────────────────────────────────────────────────────────────
\usepackage[colorlinks=true, linkcolor=azulNoche, urlcolor=cyanNeon]{hyperref}
\sloppy
\emergencystretch 3em

% ── Paleta de colores — tecnológico dark-mode ──────────────────────────────────
\definecolor{fondoOscuro}{HTML}{0D1B2A}%             Dark navy — fondo banda título
\definecolor{verdeTurquesa}{HTML}{004D40}%           Verde turquesa — bandas de marca
\definecolor{azulNoche}{HTML}{1B3A6B}%               Midnight blue — secciones, marcos
\definecolor{azulMedio}{HTML}{2A5298}%               Azul medio — variante de acento
\definecolor{cyanNeon}{HTML}{00C8E0}%                Cyan eléctrico — acentos primarios
\definecolor{verdeSignal}{HTML}{00B887}%             Verde señal — fortalezas, éxito
\definecolor{naranjaVivo}{HTML}{FF7043}%             Naranja vivo — mejoras, advertencias
\definecolor{rojoAlerta}{HTML}{E53935}%              Rojo alerta — errores
\definecolor{grisPapel}{HTML}{F4F6FA}%               Fondo tarjetas claras
\definecolor{grisLinea}{HTML}{C8D0E0}%               Bordes sutiles
\definecolor{grisTexto}{HTML}{6B7A99}%               Texto secundario
\definecolor{amarilloNota}{HTML}{FFB300}%            Notas / advertencias doradas
\definecolor{codigoFondo}{HTML}{1E2D3D}%             Fondo del bloque de código
\definecolor{codigoComentario}{HTML}{7A8FA6}%        Comentarios
\definecolor{codigoCadena}{HTML}{FF9D5C}%            Cadenas / strings
\definecolor{codigoKeyword}{HTML}{00C8E0}%           Palabras clave
\definecolor{codigoNumero}{HTML}{00E5A0}%            Números

% ── Banda de título: portada infográfica ──────────────────────────────────────
%   Diseño en 2 capas: banda de color (por defecto fondo oscuro) + franja de
%   acento cyan abajo. Uso: \bandaTitulo[color]{título}{subtítulo}.
%   El icono se puede cambiar con \renewcommand{\iconoBanda}{...} (FontAwesome).
\newcommand{\iconoBanda}{microchip}
\newcommand{\bandaTitulo}[3][fondoOscuro]{%
  \begin{tcolorbox}[
    enhanced,
    colback=#1, colframe=#1,
    boxrule=0pt, arc=8pt, outer arc=8pt,
    width=\linewidth,
    top=18pt, bottom=6pt, left=14pt, right=14pt,
    borderline south={3.5pt}{0pt}{cyanNeon},
  ]
    \begin{minipage}[c]{0.07\linewidth}
      \centering
      \textcolor{cyanNeon}{\fontsize{30}{30}\selectfont\faIcon{\iconoBanda}}
    \end{minipage}%
    \hspace{8pt}%
    \begin{minipage}[c]{0.88\linewidth}
      {\fontsize{20}{24}\selectfont\bfseries\color{white}#2}\par\vspace{5pt}%
      \textcolor{cyanNeon!80!white}{\small\faIcon{angle-right}~#3}%
    \end{minipage}
    \vspace{4pt}
  \end{tcolorbox}%
  \vspace{8pt}%
}

% ── Tarjetas de datos (infografía) ────────────────────────────────────────────
\tcbset{
  tarjetaDato/.style={
    enhanced, breakable,
    arc=6pt, outer arc=6pt,
    colback=grisPapel, colframe=azulNoche,
    boxrule=1pt,
    fonttitle=\bfseries\small, coltitle=white,
    attach boxed title to top left={yshift=-3mm, xshift=6mm},
    boxed title style={
      colback=azulNoche, colframe=azulNoche,
      arc=4pt, boxrule=0pt,
    },
    drop fuzzy shadow=azulNoche!30!white,
    top=8pt, bottom=8pt, left=10pt, right=10pt,
  },
  tarjetaIcono/.style={
    enhanced, breakable,
    arc=6pt, outer arc=6pt,
    colback=white, colframe=grisLinea, boxrule=0.8pt,
    fonttitle=\bfseries\small, coltitle=azulNoche,
    borderline west={3pt}{0pt}{cyanNeon},
    drop fuzzy shadow=azulNoche!25!white,
    top=6pt, bottom=6pt, left=10pt, right=10pt,
  },
  cajaTitulo/.style={
    enhanced, breakable,
    arc=6pt, outer arc=6pt,
    colback=azulNoche!10!grisPapel, colframe=azulNoche,
    boxrule=1pt,
    fonttitle=\bfseries, coltitle=white,
    attach boxed title to top left={yshift=-3mm, xshift=6mm},
    boxed title style={
      colback=azulNoche, colframe=azulNoche,
      arc=4pt, boxrule=0pt,
    },
    drop fuzzy shadow=azulNoche!25!white,
    top=8pt, bottom=8pt, left=10pt, right=10pt,
  },
  cajaContenido/.style={
    enhanced, breakable,
    arc=5pt, outer arc=5pt,
    colback=grisPapel, colframe=grisLinea,
    boxrule=0.8pt,
    fonttitle=\bfseries\small, coltitle=azulNoche,
    top=6pt, bottom=6pt, left=10pt, right=10pt,
  },
  cajaFortaleza/.style={
    enhanced, breakable,
    arc=6pt, outer arc=6pt,
    colback=verdeSignal!8!white, colframe=verdeSignal,
    boxrule=1pt,
    borderline west={4pt}{0pt}{verdeSignal},
    fonttitle=\bfseries\small, coltitle=verdeSignal!70!black,
    drop fuzzy shadow=verdeSignal!20!white,
    top=6pt, bottom=6pt, left=10pt, right=10pt,
  },
  cajaMejora/.style={
    enhanced, breakable,
    arc=6pt, outer arc=6pt,
    colback=naranjaVivo!8!white, colframe=naranjaVivo,
    boxrule=1pt,
    borderline west={4pt}{0pt}{naranjaVivo},
    fonttitle=\bfseries\small, coltitle=naranjaVivo!80!black,
    drop fuzzy shadow=naranjaVivo!20!white,
    top=6pt, bottom=6pt, left=10pt, right=10pt,
  },
  cajaRecomendacion/.style={
    enhanced, breakable,
    arc=6pt, outer arc=6pt,
    colback=cyanNeon!6!white, colframe=cyanNeon!60!azulNoche,
    boxrule=1pt,
    borderline west={4pt}{0pt}{cyanNeon},
    fonttitle=\bfseries\small, coltitle=azulNoche,
    drop fuzzy shadow=azulNoche!20!white,
    top=6pt, bottom=6pt, left=10pt, right=10pt,
  },
}

% ── Ícono + texto (encabezados de sección y elementos) ────────────────────────
\newcommand{\iconotexto}[2]{\textcolor{cyanNeon}{\faIcon{#1}}~#2}

% ── Listados de código: estilo dark + auto-envoltura ─────────────────────────
\lstdefinestyle{estiloCodigo}{
    backgroundcolor=\color{codigoFondo},
    basicstyle=\ttfamily\small\color{white},
    commentstyle=\color{codigoComentario}\itshape,
    stringstyle=\color{codigoCadena},
    keywordstyle=\color{codigoKeyword}\bfseries,
    numberstyle=\tiny\color{codigoComentario},
    numbers=left,
    numbersep=10pt,
    showstringspaces=false,
    breaklines=true,
    breakatwhitespace=false,
    tabsize=4,
    frame=none,
    captionpos=b,
    aboveskip=6pt,
    belowskip=4pt,
    xleftmargin=14pt,
    framexleftmargin=14pt,
    literate=
      {á}{{\'a}}1 {é}{{\'e}}1 {í}{{\'i}}1 {ó}{{\'o}}1 {ú}{{\'u}}1
      {Á}{{\'A}}1 {É}{{\'E}}1 {Í}{{\'I}}1 {Ó}{{\'O}}1 {Ú}{{\'U}}1
      {ñ}{{\~n}}1 {Ñ}{{\~N}}1 {ü}{{\"u}}1 {Ü}{{\"U}}1
      {¿}{{?`}}1 {¡}{{!`}}1
}
\lstdefinestyle{estiloPython}{
    style=estiloCodigo,
    language=Python,
}
\lstset{style=estiloCodigo}
\BeforeBeginEnvironment{lstlisting}{%
  \begin{tcolorbox}[
    enhanced, breakable,
    arc=6pt, outer arc=6pt,
    colback=codigoFondo, colframe=azulNoche,
    boxrule=1pt,
    drop fuzzy shadow=azulNoche!25!white,
    top=2pt, bottom=2pt, left=0pt, right=0pt,
  ]%
}
\AfterEndEnvironment{lstlisting}{\end{tcolorbox}}

% ── Cajas de contenido temáticas ──────────────────────────────────────────────
\newtcolorbox{cajaRecuerda}{
    enhanced, breakable,
    arc=6pt, outer arc=6pt,
    colback=azulNoche!10!grisPapel,
    colframe=azulNoche, boxrule=1pt,
    borderline west={4pt}{0pt}{cyanNeon},
    fonttitle=\bfseries\small\color{white},
    title={\faIcon{info-circle}~Recuerda},
    attach boxed title to top left={yshift=-3mm, xshift=6mm},
    boxed title style={colback=azulNoche, arc=4pt, boxrule=0pt},
    drop fuzzy shadow=azulNoche!25!white,
    left=10pt, right=10pt, top=8pt, bottom=8pt,
}

\newtcolorbox{cajaConcepto}{
    enhanced, breakable,
    arc=6pt, outer arc=6pt,
    colback=verdeSignal!8!white,
    colframe=verdeSignal, boxrule=1pt,
    borderline west={4pt}{0pt}{verdeSignal},
    fonttitle=\bfseries\small\color{white},
    title={\faIcon{lightbulb}~Concepto clave},
    attach boxed title to top left={yshift=-3mm, xshift=6mm},
    boxed title style={colback=verdeSignal!80!black, arc=4pt, boxrule=0pt},
    drop fuzzy shadow=verdeSignal!20!white,
    left=10pt, right=10pt, top=8pt, bottom=8pt,
}

\newtcolorbox{cajaEjemplo}{
    enhanced, breakable,
    arc=6pt, outer arc=6pt,
    colback=cyanNeon!5!white,
    colframe=cyanNeon!70!azulNoche, boxrule=1pt,
    borderline west={4pt}{0pt}{cyanNeon},
    fonttitle=\bfseries\small\color{fondoOscuro},
    title={\faIcon{code}~Ejemplo},
    attach boxed title to top left={yshift=-3mm, xshift=6mm},
    boxed title style={colback=cyanNeon, arc=4pt, boxrule=0pt},
    drop fuzzy shadow=azulNoche!20!white,
    left=10pt, right=10pt, top=8pt, bottom=8pt,
}

\newtcolorbox{cajaEjercicio}{
    enhanced, breakable,
    arc=6pt, outer arc=6pt,
    colback=azulMedio!7!white,
    colframe=azulMedio, boxrule=1pt,
    borderline west={4pt}{0pt}{azulMedio},
    fonttitle=\bfseries\small\color{white},
    title={\faIcon{pencil-alt}~Ejercicio},
    attach boxed title to top left={yshift=-3mm, xshift=6mm},
    boxed title style={colback=azulMedio, arc=4pt, boxrule=0pt},
    drop fuzzy shadow=azulNoche!20!white,
    left=10pt, right=10pt, top=8pt, bottom=8pt,
}

\newtcolorbox{cajaReto}{
    enhanced, breakable,
    arc=6pt, outer arc=6pt,
    colback=naranjaVivo!6!white,
    colframe=naranjaVivo, boxrule=1pt,
    borderline west={4pt}{0pt}{naranjaVivo},
    fonttitle=\bfseries\small\color{white},
    title={\faIcon{fire}~Reto},
    attach boxed title to top left={yshift=-3mm, xshift=6mm},
    boxed title style={colback=naranjaVivo, arc=4pt, boxrule=0pt},
    drop fuzzy shadow=naranjaVivo!20!white,
    left=10pt, right=10pt, top=8pt, bottom=8pt,
}

% ── Indicadores de dificultad ─────────────────────────────────────────────────
\newcommand{\facil}{\textcolor{verdeSignal}{\faIcon{circle}\,\textbf{Fácil}}}
\newcommand{\intermedio}{\textcolor{naranjaVivo}{\faIcon{circle}\,\textbf{Intermedio}}}
\newcommand{\dificil}{\textcolor{rojoAlerta}{\faIcon{circle}\,\textbf{Difícil}}}

% ─────────────────────────────────────────────────────────────────────────────
\renewcommand{\iconoBanda}{flask}
\fancyhead[L]{\footnotesize\color{grisTexto}\faIcon{flask}~LABORATORIO\_I\_FISICA}
\fancyhead[R]{\footnotesize\color{grisTexto}Práctica 08}
\hypersetup{pdftitle={Laboratorio I de Física},pdfauthor={Prof. César Rodríguez}
  pdfsubject={LABORATORIO\_I\_FISICA}}

\begin{document}
\bandaTitulo[fondoOscuro]{Práctica 08 --- Ley de Hooke y constante elástica}{Semana 8 --- Fase 3: Elasticidad, oscilaciones y calor}

\begin{center}\small\color{grisTexto}Prof. César Rodríguez \textbar{} Universidad de Oriente, Núcleo de Sucre, Departamento de Física \textbar{} 6 de noviembre de 2026\end{center}

\begin{tabularx}{\linewidth}{@{}>{\raggedright\arraybackslash}p{4.6cm}X@{}}
\toprule
\textbf{Duración} & 3 horas de laboratorio presencial más 2 horas de análisis e informe \\
\textbf{Modalidad} & 4 grupos con roles rotativos \\
\textbf{Entrega} & GitHub del grupo, carpeta semana\_08, antes del corte semanal \\
\textbf{Instrumento} & Práctica con informe \\
\bottomrule
\end{tabularx}

\section{\iconotexto{bullseye}{Objetivo}}

Verificar la proporcionalidad entre la fuerza aplicada y el estiramiento de un resorte, y determinar su constante elástica a partir de la pendiente de un ajuste por mínimos cuadrados.

\section{\iconotexto{clipboard-list}{Contenidos que se aplican}}

\begin{itemize}[leftmargin=1.4em,itemsep=2pt,topsep=3pt]
  \item Ley de Hooke y régimen elástico.
  \item Fuerza aplicada como peso de masas calibradas.
  \item Estiramiento respecto de la posición de equilibrio sin carga.
  \item Ajuste lineal de fuerza frente a estiramiento.
  \item Constante elástica, unidades y significado físico.
  \item Límite elástico y desviación de la linealidad.
\end{itemize}

\section{\iconotexto{book-open}{Fundamento teórico}}

Dentro de su régimen elástico, un resorte ejerce una fuerza recuperadora
proporcional a su deformación:
\[
  F = -k\,\Delta y.
\]
Si se cuelga una masa $m$ y el sistema queda en equilibrio, la fuerza
aplicada es su peso y
\[
  m\,g = k\,\Delta y
  \qquad\Longrightarrow\qquad
  k = \frac{m\,g}{\Delta y}.
\]
El estiramiento $\Delta y$ se mide respecto de la posición del resorte sin
carga, no respecto de su longitud natural si el propio porta-pesas ya lo
estira. Al graficar $P = m\,g$ frente a $\Delta y$ se obtiene una recta
cuya pendiente es $k$:
\[
  P = k\,\Delta y.
\]
La constante se toma de la pendiente del ajuste, no del promedio de los
cocientes $P_i/\Delta y_i$, porque el ajuste pondera todos los puntos y
reduce el efecto del ruido de lectura. La incertidumbre de $k$ se obtiene
del ajuste. Si los residuales se curvan en los extremos, el resorte salió
del régimen elástico o las espiras se tocaron.

\begin{center}
\begin{tikzpicture}[scale=0.9]
  \draw[thick] (0.5,6.2) -- (4.5,6.2);
  \draw[thick] (0.5,6.2) -- (0.5,5.9);
  \draw[thick] (4.5,6.2) -- (4.5,5.9);
  \draw[thick] (1.2,5.9) -- (3.8,5.9);
  \draw[thick] (2.5,5.9) -- (2.5,5.5);
  \draw[thick] (2.5,5.5)
    -- (2.2,5.3) -- (2.8,5.1) -- (2.2,4.9) -- (2.8,4.7)
    -- (2.2,4.5) -- (2.8,4.3) -- (2.2,4.1) -- (2.8,3.9)
    -- (2.2,3.7) -- (2.8,3.5) -- (2.2,3.3) -- (2.8,3.1)
    -- (2.5,2.9);
  \draw[thick] (2.5,2.9) -- (2.5,2.6);
  \fill[gray] (2.0,1.6) rectangle (3.0,2.6);
  \draw[dashed] (0.9,5.7) -- (4.1,5.7);
  \draw[dashed] (0.9,2.6) -- (4.1,2.6);
  \draw[<->] (3.7,5.7) -- (3.7,2.6);
  \node[right] at (3.75,4.15) {$\Delta y$};
\end{tikzpicture}
\par\smallskip{\small\itshape Resorte con masas colgadas: el estiramiento se mide desde la posición sin carga}
\end{center}

\section{\iconotexto{box-open}{Materiales y equipos}}

\begin{itemize}[leftmargin=1.4em,itemsep=2pt,topsep=3pt]
  \item Soporte universal con pinza y nuez
  \item Resorte helicoidal de constante desconocida
  \item Juego de masas de 50 g a 500 g
  \item Regla milimetrada de 50 cm
  \item Balanza digital de 0,1 g
  \item Cronómetro para el método dinámico opcional
\end{itemize}

\section{\iconotexto{flask}{Escenario}}

Cada grupo cuelga masas crecientes de un resorte, mide el estiramiento y ajusta la recta fuerza-estiramiento. Debe reportar la constante elástica con su incertidumbre y señalar hasta qué masa el resorte sigue siendo lineal.

\section{\iconotexto{list-ol}{Procedimiento}}

\begin{enumerate}[leftmargin=1.4em,itemsep=2pt,topsep=3pt]
  \item Montar el soporte con el resorte vertical y una regla fija junto a él, sin tocar el resorte.
  \item Anotar la posición del extremo inferior del resorte sin carga; esa es la referencia para el estiramiento.
  \item Colgar masas crecientes de 50 g en 50 g hasta 500 g, esperando a que el sistema se detenga antes de leer.
  \item Anotar la posición del extremo inferior para cada masa y calcular el estiramiento respecto de la referencia.
  \item Retirar las masas una a una y comprobar que el resorte recupera su longitud inicial; si no, se superó el límite elástico.
  \item Convertir cada masa a peso y ajustar la recta por mínimos cuadrados.
  \item Reportar la constante elástica con su incertidumbre y el coeficiente de determinación.
\end{enumerate}

\section{\iconotexto{table}{Tabla de datos}}

\begin{center}
\begin{minipage}{\linewidth}
\centering
{\small
\textbf{Masa colgada, peso y estiramiento del resorte}\par\smallskip
\begin{tabular}{@{}ccccc@{}}
\toprule
Ensayo & $m$ (g) & $P$ (N) & $y_i$ (cm) & $\Delta y$ (cm) \\
\midrule
1 & 50 & $\ldots$ & $\ldots$ & $\ldots$ \\
2 & 100 & $\ldots$ & $\ldots$ & $\ldots$ \\
3 & 150 & $\ldots$ & $\ldots$ & $\ldots$ \\
4 & 200 & $\ldots$ & $\ldots$ & $\ldots$ \\
5 & 250 & $\ldots$ & $\ldots$ & $\ldots$ \\
6 & 300 & $\ldots$ & $\ldots$ & $\ldots$ \\
7 & 350 & $\ldots$ & $\ldots$ & $\ldots$ \\
8 & 400 & $\ldots$ & $\ldots$ & $\ldots$ \\
9 & 450 & $\ldots$ & $\ldots$ & $\ldots$ \\
10 & 500 & $\ldots$ & $\ldots$ & $\ldots$ \\
\bottomrule
\end{tabular}
}\par\smallskip{\footnotesize El estiramiento se mide desde la posición sin carga; si se usa el método dinámico, anotar aparte el periodo de diez oscilaciones para cada masa.}
\end{minipage}
\end{center}

\section{\iconotexto{calculator}{Análisis y tratamiento de datos}}

\begin{enumerate}[leftmargin=1.4em,itemsep=2pt,topsep=3pt]
  \item Calcular el peso de cada masa y el estiramiento respecto de la posición sin carga.
  \item Graficar el peso frente al estiramiento con barras de incertidumbre.
  \item Ajustar la recta por mínimos cuadrados y reportar la pendiente con sus unidades.
  \item Calcular la constante elástica en N/m y su incertidumbre a partir del ajuste.
  \item Analizar los residuales y decidir en qué rango el resorte cumple la ley de Hooke.
  \item Comparar el valor estático con el que se obtendría del método dinámico, si se realizó, y discutir la diferencia.
\end{enumerate}

\section{\iconotexto{plus-circle}{Ampliación (opcional)}}

\begin{itemize}[leftmargin=1.4em,itemsep=2pt,topsep=3pt]
  \item Método dinámico: colgar cada masa, separarla ligeramente de la vertical y medir el periodo de diez oscilaciones. Al graficar $T^2$ frente a $m$ se obtiene una recta cuya pendiente es $4\pi^2/k$. Compare la constante dinámica con la estática y discuta las fuentes de diferencia.
\end{itemize}

\section{\iconotexto{box-open}{Entregable}}

Informe con la tabla de datos, la gráfica con el ajuste y los residuales, la constante elástica con su incertidumbre y el rango de validez de la ley de Hooke.

\section{\iconotexto{ruler-combined}{Rúbrica de corrección}}

\begin{tabularx}{\linewidth}{@{}X r@{}}
\toprule
Criterio & Peso\\
\midrule
Fundamento teórico y modelo aplicado & 1,0 \\
Montaje, seguridad y manejo del instrumental & 1,5 \\
Toma de datos completa y trazable & 2,0 \\
Presupuesto de incertidumbres & 2,0 \\
Análisis, ajuste y gráficas & 2,0 \\
Conclusiones contrastadas con el modelo & 1,0 \\
Presentación y estructura del informe & 0,5 \\
\midrule
\textbf{Total} & \textbf{10,0}\\
\bottomrule
\end{tabularx}

\begin{tcolorbox}[cajaMejora]\textbf{Antes de empezar:} revisa el montaje contra la figura y verifica el cero y el rango de cada instrumento antes de tomar la primera lectura. Si algo no coincide, detente y consulta al docente: los datos tomados sobre un montaje dudoso no se corrigen después.
\end{tcolorbox}
\end{document}
